% MT038 / explanation — canonical editable content.
题目给定两者准确率和总计算量相同，比较的是100核芯片上一次\textbf{预测的等待时间}。总工作量相同，不代表能够并行完成的程度相同。


RF有10棵深度10的树。不同树可以同时预测，但单棵树必须先判断根节点，才能知道下一步访问哪个分支；同一路径上的约10次判断前后相依。简单地“一棵树一个核心”时，只有10个主要并行任务。


若AdaBoost的100个弱学习器都很浅，例如每个只是一个树桩，那么训练完成后，它们对同一输入的预测可同时放到100个核心上。最后把100个结果加权求和，再输出类别。弱学习器的\textbf{训练}需要前一轮的样本权重，但这不要求它们的\textbf{预测}依次执行。


理想化地，若一个树桩一次判断，RF路径需10次，每次耗时相同且忽略汇总与调度，AdaBoost的主体预测可比RF快约10倍。这是说明并行度的算例，不是由算法名称保证的真实倍率。


因此按原题意选择AdaBoost，同时交代“弱学习器浅、能够并行、汇总开销较小”的假设。如果100个弱学习器很复杂，或核心通信开销大，仅给数量就不足以推出相同优势。


\textbf{出处：} 2021B*期中 Q12，原卷第4页（10分）


\textbf{答案依据：} 2021B*期中 Q12 随卷／配套参考答案；条件与解释经复核。
